% ============================================================
% Appendix A: Supplementary Material
% ============================================================
\chapter{Supplementary Material}
\label{ch:appendix-a}

\section{Proof of \cref{thm:regret}}
\label{sec:proof-main}

We provide the complete proof of the main regret bound.

\begin{proofbox}[Proof of \cref{thm:regret}]
The proof proceeds in three steps.

\textbf{Step 1: Decomposition of the regret.}
The simple regret can be decomposed as:
\begin{equation}
    r_n = f(\bx^\star) - f(\bx_n^\star) \leq
    \alpha_n(\bx_n^\star) + \varepsilon_n(\bx_n^\star),
    \label{eq:regret-decompose}
\end{equation}
where $\alpha_n(\bx_n^\star)$ is the value of the acquisition function at the
chosen point and $\varepsilon_n(\bx_n^\star)$ captures the approximation error
due to the adaptive subspace.

\textbf{Step 2: Bounding the acquisition function.}
Using the standard \gls{gp} regret bound and the adaptive subspace structure:
\begin{equation}
    \alpha_n(\bx) \leq \sqrt{2\kappa_{\max}^2 \beta_n} \cdot \sigma_n(\bx),
    \label{eq:acquisition-bound}
\end{equation}
where $\sigma_n(\bx)$ is the posterior standard deviation in the projected
space and $\beta_n$ is the exploration parameter.

\textbf{Step 3: Summing over the subspace.}
By \cref{lem:subspace}, the effective dimensionality of the projected space is
$k_n$, and the maximum information gain satisfies $\gamma_n \leq
k_n \log(1 + n^2 \pi^2 / (6\delta))$.  Combining these bounds yields:
\begin{equation}
    r_n \leq \sqrt{2\kappa_{\max}^2 \beta_n \cdot \frac{k_n \gamma_n}{n}}.
    \label{eq:final-bound}
\end{equation}
\end{proofbox}

\section{Proof of \cref{lem:subspace}}
\label{sec:proof-subspace}

\begin{proofbox}[Proof of \cref{lem:subspace}]
We bound the Frobenius norm error $\| \bZ_n - \bZ^\star \|_F$ using the
variational lower bound.  By the \gls{kl} divergence inequality:
\begin{equation}
    \KL(q(\bZ_n) \| p(\bZ \mid \cD_n)) \leq \frac{1}{2} \Tr(\bSigma_n^{-1}
    \bSigma^\star) + \frac{1}{2} \| \bmu_n - \bmu^\star \|_{\bSigma_n^{-1}}^2
    - \frac{k d}{2},
    \label{eq:kl-expansion}
\end{equation}
where $\bmu_n, \bSigma_n$ are the variational parameters and
$\bmu^\star, \bSigma^\star$ are the true posterior parameters.

Since $\KL(q \| p) \leq C / \sqrt{n}$ by the variational inference convergence
result (Theorem 1 of~\citet{bibid}), we obtain:
\begin{equation}
    \| \bZ_n - \bZ^\star \|_F \leq \sqrt{\frac{2C}{\sqrt{n}}} = \frac{C'}{\sqrt[4]{n}},
    \label{eq:frobenius-bound}
\end{equation}
where $C' = \sqrt{2C}$ depends on the prior parameters and the smoothness
of $f$.
\end{proofbox}

\section{Extended Results Tables}
\label{sec:extended-results}

\Cref{tab:extended-d20} and \cref{tab:extended-d500} provide the complete
results at $d=20$ and $d=500$ respectively.

\begin{table}[htbp]
    \centering
    \caption{Final simple regret at $d=20$ ($n=500$ evaluations).}
    \label{tab:extended-d20}
    \begin{tabularx}{\textwidth}{@{}lXXXXX@{}}
        \toprule
        \textbf{Function} & \textbf{ASBO} & \textbf{REMBO} & \textbf{TuRBO} &
        \textbf{PESBO} & \textbf{SAASBO} \\
        \midrule
        Branin-Hartmann & $\mathbf{2.1 \times 10^{-4}}$ & $5.8 \times 10^{-4}$ &
        $3.2 \times 10^{-4}$ & $2.8 \times 10^{-4}$ & $4.5 \times 10^{-4}$ \\
        Ackley & $\mathbf{3.5 \times 10^{-4}}$ & $8.2 \times 10^{-4}$ &
        $5.1 \times 10^{-4}$ & $4.2 \times 10^{-4}$ & $7.1 \times 10^{-4}$ \\
        Rosenbrock & $\mathbf{1.8 \times 10^{-4}}$ & $4.1 \times 10^{-4}$ &
        $2.5 \times 10^{-4}$ & $2.1 \times 10^{-4}$ & $3.8 \times 10^{-4}$ \\
        Schwefel & $\mathbf{2.9 \times 10^{-4}}$ & $7.5 \times 10^{-4}$ &
        $4.8 \times 10^{-4}$ & $3.6 \times 10^{-4}$ & $6.2 \times 10^{-4}$ \\
        Rastrigin & $\mathbf{1.5 \times 10^{-4}}$ & $3.9 \times 10^{-4}$ &
        $2.8 \times 10^{-4}$ & $1.9 \times 10^{-4}$ & $3.3 \times 10^{-4}$ \\
        \bottomrule
    \end{tabularx}
\end{table}

\begin{table}[htbp]
    \centering
    \caption{Final simple regret at $d=500$ ($n=500$ evaluations).}
    \label{tab:extended-d500}
    \begin{tabularx}{\textwidth}{@{}lXXXXX@{}}
        \toprule
        \textbf{Function} & \textbf{ASBO} & \textbf{REMBO} & \textbf{TuRBO} &
        \textbf{PESBO} & \textbf{SAASBO} \\
        \midrule
        Branin-Hartmann & $\mathbf{4.2 \times 10^{-3}}$ & $1.8 \times 10^{-2}$ &
        $3.1 \times 10^{-2}$ & $8.5 \times 10^{-3}$ & $2.1 \times 10^{-2}$ \\
        Ackley & $\mathbf{8.1 \times 10^{-3}}$ & $3.5 \times 10^{-2}$ &
        $5.2 \times 10^{-2}$ & $1.4 \times 10^{-2}$ & $4.2 \times 10^{-2}$ \\
        Rosenbrock & $\mathbf{6.5 \times 10^{-3}}$ & $2.8 \times 10^{-2}$ &
        $4.1 \times 10^{-2}$ & $1.1 \times 10^{-2}$ & $3.3 \times 10^{-2}$ \\
        Schwefel & $\mathbf{5.8 \times 10^{-3}}$ & $2.4 \times 10^{-2}$ &
        $3.8 \times 10^{-2}$ & $9.2 \times 10^{-3}$ & $2.9 \times 10^{-2}$ \\
        Rastrigin & $\mathbf{3.9 \times 10^{-3}}$ & $1.6 \times 10^{-2}$ &
        $2.7 \times 10^{-2}$ & $6.8 \times 10^{-3}$ & $1.9 \times 10^{-2}$ \\
        \bottomrule
    \end{tabularx}
\end{table}

\section{Reproducibility Checklist}
\label{sec:reproducibility}

To ensure reproducibility, we provide the following checklist:

\begin{enumerate}
    \item \textbf{Code availability:} The complete source code for ASBO and all
    experiments is available at
    \url{https://github.com/example/highdimbo} under the MIT licence.

    \item \textbf{Random seeds:} All experiments use fixed random seeds
    (\texttt{seed=42} for each of the 20 independent runs).

    \item \textbf{Hyperparameters:} The complete hyperparameter configuration
    is provided in \texttt{configs/} directory of the repository.

    \item \textbf{Dependencies:} A \texttt{requirements.txt} file specifies
    exact package versions for reproducibility.

    \item \textbf{Computational requirements:} Full replication requires
    approximately 48 GPU-hours on an NVIDIA A100.

    \item \textbf{Data:} All synthetic benchmark functions are standard and
    defined in the code.  Real-world datasets are publicly available and
    referenced in the paper.

    \item \textbf{License:} The code is released under the MIT licence, and
    the data under CC-BY-4.0.
\end{enumerate}

\section{Hyperparameter Sensitivity}
\label{sec:hyperparameter-sensitivity}

We conducted a sensitivity analysis of the key hyperparameters of ASBO.
\Cref{tab:hyperparameters} summarises the default values and the range over
which each parameter was varied.

\begin{table}[htbp]
    \centering
    \caption{Hyperparameter sensitivity analysis.  Default values and tested
    ranges.}
    \label{tab:hyperparameters}
    \begin{tabularx}{\textwidth}{@{}lXcc@{}}
        \toprule
        \textbf{Parameter} & \textbf{Description} & \textbf{Default} &
        \textbf{Range} \\
        \midrule
        $k_0$ & Initial subspace dimension & 10 & $\{5, 10, 20, 50\}$ \\
        $n_0$ & Initial samples & $2(d+1)$ & $\{d, 2d, 4d\}$ \\
        $\kappa$ & Subspace update frequency & 5 & $\{1, 5, 10, 20\}$ \\
        $\delta$ & Variational inference precision & $0.01$ & $\{0.001, 0.01, 0.1\}$ \\
        $\ell$ & \gls{gp} length-scale & $1.0$ & $\{0.1, 1.0, 10.0\}$ \\
        $\sigma_f$ & Signal variance & $1.0$ & $\{0.1, 1.0, 5.0\}$ \\
        \bottomrule
    \end{tabularx}
\end{table}
